\section{Related Work}
\label{sec:related}

\subsection{DRL for Mapless Navigation}

Tai et al.~\cite{tai2017virtual} demonstrated that a policy trained
entirely in simulation can transfer to physical robots using raw
2D LiDAR observations, establishing the viability of end-to-end
sim-to-real navigation without explicit mapping. Pfeiffer et
al.~\cite{pfeiffer2017dl} extended this paradigm to data-driven motion
planning, showing strong generalization across varied indoor layouts.
These approaches typically pair a 1D convolutional feature extractor
with an off-policy actor-critic algorithm such as
TD3~\cite{fujimoto2018td3}, the architecture we adopt as our baseline.
While effective in randomized obstacle environments, such policies are
evaluated almost exclusively on metrics aggregated over diverse but
unstructured scenes, which we show in Section~\ref{sec:results} masks
systematic failure modes in concave or symmetric geometries.

\subsection{Memory-Augmented Navigation}

Recurrent architectures have been proposed to give navigation policies
access to temporal context beyond the instantaneous sensor reading,
under the assumption that memory of recent observations aids decision-
making in partially observable settings. We test this hypothesis with
a GRU-based variant (RCPG) and find the effect to be scenario-
dependent rather than uniformly beneficial: temporal memory improves
performance in tasks requiring precise alignment (narrow doorways) and
breaks deadlocks caused by bilateral symmetry, but actively harms
performance in concave trap geometries by reinforcing path persistence
when backtracking is required. This double-edged effect, to our
knowledge, has not been systematically characterized in prior
mapless-navigation literature.

\subsection{Local Minima and Trap Escape}

Classical reactive methods such as the Bug family of algorithms
guarantee trap escape through deterministic wall-following but require
explicit boundary-tracing logic. Within the learning-based literature,
recent work addresses concave-obstacle failure through auxiliary
mechanisms: hybrid potential-field and wall-following controllers
require an explicit trap detector, while environment-prediction
approaches rely on learned generative models of the surrounding
geometry. In contrast, our approach requires no additional inference
module at deployment time -- the trap-escape behavior is encoded
entirely through the training-time reward and curriculum, leaving the
network architecture and runtime inference cost unchanged from the
baseline.

\subsection{Model-Based Local Planning: NeuPAN}

NeuPAN~\cite{han2025neupan} represents the current state of the art in
model-based, real-time, map-free motion planning, directly mapping raw
LiDAR points to control actions through a proximal alternating-
minimization network without intermediate occupancy-grid construction.
NeuPAN reports strong performance across corridor, parking-lot, and
cluttered-sandbox benchmarks using robot platforms on the order of
1.6\,m in length. We evaluate NeuPAN under two configurations against
our structured benchmark: (1) its original robot geometry, and (2) a
DUNE perception model retrained for TurtleBot3-scale dimensions
(0.178$\times$0.138\,m). We find that even with correctly-matched
robot parameters, NeuPAN's local optimization can become trapped in a
persistent oscillatory local minimum near constrained multi-obstacle
geometries -- a failure mode distinct from, and complementary to, its
previously documented limitations with fast-moving dynamic obstacles.
This suggests that constrained corridor navigation remains challenging
for purely local, optimization-based planners regardless of parameter
tuning, motivating the learning-based curriculum approach explored in
this work.
